\section{Owner-count and resource-dimension extensions}\label{ext:sec}
The hybrid of Section~\ref{prod:section} extends to every finite monotone,
concave, degree-one homogeneous production function, including discontinuities
at boundary faces.  Its budget can also count the selected forms rather than
the entire technology.  We obtain an $n/\log n$ owner-count upper bound,
an imported $\log n/\log\log n$ lower bound, and a dimension fallback that is
sharp without concavity.  The sharp rates in owners and in concave resource
dimension remain open.

Throughout, $f:\R_+^d\to\R$ is finite, coordinatewise nondecreasing and
positively homogeneous of degree one; concavity is assumed where stated.
Homogeneity forces $f(0)=0$, hence monotonicity gives $f\ge0$.
There are $n$ fixed owners $N=[n]$, with arbitrary bundles $b_i\ge0$;
$c=\sum_i b_i$, $T=f(c)$, $v_b(S)=f(\sum_{i\in S}b_i)$, and
$x_B=\sum_iB_i b_i$.  An eligible rule is nonnegative, coalition-covering,
and componentwise nondecreasing in each owner's whole report at every real
profile, as in Section~\ref{prod:section}.  All domains and owners remain fixed.
Use $\alpha_j,C_j(H)$ from \eqref{prod:constant}, with $H\ge1$ in place of $K$.

\subsection{Reusable estimates and the selected range}
\begin{lemma}[Fixed-vector Chernoff bound]\label{ext:lem:chernoff}
Let independent bits have $\Prob(B_i=0)=q\in(0,1)$ and put
$\alpha=-\log(q+(1-q)e^{-\lambda})$ for $\lambda>0$.
For every fixed $\delta\in[0,1]^n$,
\[
 \Prob\Bigl(\sum_iB_i\delta_i\le1\Bigr)
 \le e^{\lambda-\alpha\sum_i\delta_i}.
\]
If $\mathcal L\subseteq[0,1]^n$ has at most $H$ elements and a mask-dependent
choice $\delta^B\in\mathcal L$ always satisfies
$\sum_iB_i\delta_i^B\le1$, then
$\E\sum_i\delta_i^B\le(\log H+\lambda+1)/\alpha$.
For $q\le1/j$, $\lambda=2\log j$, one has $\alpha\ge\alpha_j$.
\end{lemma}
\begin{proof}
The fixed-vector exponential calculation in \eqref{prod:chernoff} uses only
$\delta_i\in[0,1]$ and independence, so applies verbatim.
The event $\sum_i\delta_i^B\ge x$ implies that some vector in the fixed
list with sum at least $x$ has sampled score at most one.
A union bound gives $\min\{1,He^{\lambda-\alpha x}\}$; its integral is
$(\log H+\lambda+1)/\alpha$.  Finally
$q+(1-q)/j^2\le(j^2+j-1)/j^3$ when $q\le1/j$.
The list must be fixed before sampling, but may depend on the full profile.
\end{proof}

\begin{lemma}[Budget from capped mask contributions]\label{ext:lem:budget}
Fix a profile and, for every mask, a vector $a(B)\in[0,T]^n$ satisfying
\begin{equation}\label{ext:eq:open-sum}
 \sum_iB_i a_i(B)=f(x_B).
\end{equation}
Define $z_i^q=\E a_i(B)$, $p_i^q=\E[a_i(B)\mid B_i=0]$, and use the endpoint
marginal $m_i^q$ and hybrid $r_i^q=qp_i^q+(1-q)m_i^q$ of
\eqref{prod:hybrid}.  With $g=\E f(x_B)$, $Z=\sum_i z_i^q$, $R=\sum_i r_i^q$,
all these functions are finite mask polynomials and
\[
 q\sum_i p_i^q=Z-g,\qquad \sum_i m_i^q=-g',\qquad R=Z-g-(1-q)g'.
\]
Thus $\int_0^{b_0}R\le\int_0^{b_0}Z+T$ for $0<b_0<1$.
If the vectors $a(B)$ take at most $H$ values, then for $0<q\le1/j$,
$Z(q)\le T(\log H+2\log j+1)/\alpha_j$, and
$j\int_0^{1/j}R(q)\,dq\le C_j(H)T$.
\end{lemma}
\begin{proof}
Condition on the own bit and sum \eqref{ext:eq:open-sum} to obtain the first
identity.  Differentiating the finite multiaffine expectation on its diagonal
gives the second.  The exact integration identity \eqref{prod:integral}
then applies, since $g(0)=T$ and $g\ge0$.
If $T>0$, the fixed list $\{a(B)/T\}$ lies in $[0,1]^n$, and its selected
score is $f(x_B)/T\le1$.  Apply Lemma~\ref{ext:lem:chernoff}, then integrate.
For $T=0$, monotonicity makes every masked value and contribution zero.
\end{proof}

\begin{proposition}[Selected-range refinement]\label{ext:prop:range}
For a finite min-linear family $V$ of size $K$, the unchanged hybrid rule
of Section~\ref{prod:section} satisfies
\[
 \sum_iM_i(b)\le\min\{n,C_j(\min\{K,2^n\})\}T.
\]
One may choose $j$ in advance using $\min\{K,2^n\}$; zero-report owners
are not removed and the selected range is not used to retune after reports.
\end{proposition}
\begin{proof}
At a fixed profile, $\{W(b,B):B\in\{0,1\}^n\}$ has at most
$\min\{K,2^n\}$ elements.  Its capped vectors satisfy
\eqref{ext:eq:open-sum}, because each nonnegative open contribution is a
summand of the selected value $f(x_B)\le T$ and hence is uncapped.
Apply Lemma~\ref{ext:lem:budget}.  Also $a_i,m_i^q\in[0,T]$, so
$r_i^q,M_i\le T$ and the total is at most $nT$.
\end{proof}

\subsection{General concave homogeneous production}
\begin{lemma}[Relative supporting forms]\label{ext:lem:support}
Assume $f$ is concave.  For $x\ge0$ let $J=\operatorname{supp}(x)$ and
$F_J=\{y\ge0:\operatorname{supp}(y)\subseteq J\}$.
If $J\ne\emptyset$, the set of $u\in\R_+^J$ satisfying
\[
 u\cdot x_J=f(x),\qquad f(y)\le u\cdot y_J\quad(y\in F_J)
\]
is nonempty, compact and convex.  Its unique minimum-norm element $u(x)$
therefore defines a choice depending only on $f,x$.
Set $\ell_x(y)=u(x)\cdot y_J$ on $F_J$ and $+\infty$ outside it.
For $x=0$, instead set $\ell_0(0)=0$, $\ell_0(y)=+\infty$ for $y\ne0$.
Then $\ell_x$ is componentwise nondecreasing, majorizes $f$ everywhere,
and touches it at $x$.  On its finite face it is an ordinary nonnegative
linear form.  No boundary continuity is required.
\end{lemma}
\begin{proof}
On the relative interior of $F_J$, finiteness and concavity supply a finite
supergradient $u$ at $x$~\cite{BoydSubgradients}:
$f(y)\le f(x)+u\cdot(y_J-x_J)$.
Putting $y=\lambda x$ with both $\lambda>1$ and $0<\lambda<1$ yields
$u\cdot x_J=f(x)$.  Monotonicity applied to $y=x+te_e$ gives $u_e\ge0$.
For any boundary point $y\in F_J$ and $0<\epsilon<1$, the point
$y_\epsilon=(1-\epsilon)y+\epsilon x$ is in the relative interior.
Concavity and the supporting inequality give
\[
 (1-\epsilon)f(y)+\epsilon f(x)\le f(y_\epsilon)
 \le(1-\epsilon)u\cdot y_J+\epsilon f(x).
\]
Cancel and divide by $1-\epsilon$ to obtain the boundary inequality.
The admissible set is an intersection of closed half-spaces and the touching
hyperplane.  It is bounded by $0\le u_e\le f(x)/x_e$, so is compact;
strict convexity of squared norm gives the unique choice.
The extension majorizes $f$ trivially outside its face.  Increasing an argument
within the face increases its nonnegative linear value; leaving the face
changes that value to infinity.  The zero-face convention has the same
properties.  In particular, no subtraction of infinite quantities is needed.
\end{proof}

\begin{definition}[Relative-support hybrid]\label{ext:def:general}
For each profile and mask choose the form $\ell_{x_B}$ above and set
$a_i(B)=\min\{\ell_{x_B}(b_i),T\}$.
Define $z_i^q,p_i^q,m_i^q,r_i^q$ as in Lemma~\ref{ext:lem:budget} and
$M_i=j\int_0^{1/j}r_i^q\,dq$, with $j\ge2$ fixed independently of reports.
The capped coefficients are finite, so these are finite sums and exact
integrals of polynomials, defining deterministic payments.
\end{definition}

\begin{theorem}[General concave homogeneous production]\label{ext:thm:general}
For every finite, monotone, concave, degree-one homogeneous $f$, the rule of
Definition~\ref{ext:def:general} is eligible at every real bundle profile and
\[
 0\le M_i(b)\le T,\qquad
 \sum_iM_i(b)\le\min\{n,C_j(2^n)\}T.
\]
A zero bundle receives zero, and every zero-value profile receives zero.
\end{theorem}
\begin{proof}
Both capped contributions and endpoint marginals belong to $[0,T]$, so the
same holds for their mixtures.  If $b_i=0$, then $\ell(0)=0$ and its on/off
aggregates coincide.  If $T=0$, monotonicity makes every masked value zero.

For incentives, fix $b_{-i},B_{-i}$ with $B_i=0$.
The aggregate $x_B$ and its selected form are independent of the entire own
bundle $b_i$.  Both $\ell_{x_B}(b_i)$ and $T$ increase weakly with $b_i$,
including finite-to-infinite transitions; their capped minimum is therefore
nondecreasing.  The off term of $m_i^q$ is fixed and the on term increases.
Averaging with report-independent nonnegative weights proves own-bundle
monotonicity on the whole domain.

For coverage, fix a selected form $\ell$ and coalition $S$.
If $\ell(b_i)\ge T$ for some $i\in S$, that capped term alone covers
$v_b(S)\le T$.  Otherwise all coalition terms are finite, uncapped and on
the same face, so their sum is $\ell(\sum_{i\in S}b_i)\ge v_b(S)$.
Thus every vector $a(B)$, and hence $z^q$, covers coalitions.
For a fixed $B_{-i}$ select $\ell=\ell_x$ at the on aggregate
$x=b_i+\sum_{h\ne i}B_hb_h$.  The bundles $b_i,x-b_i$ lie on its finite face.
Touching, majorization and finite additivity give
\begin{equation}\label{ext:eq:endpoint}
 f(x)-f(x-b_i)\ge\ell(x)-\ell(x-b_i)=\ell(b_i),\qquad
 0\le\ell(b_i)\le f(x)\le T.
\end{equation}
The middle term is the uncapped on contribution.  Averaging implies
$m_i^q\ge\E[a_i(B)\mid B_i=1]$, hence $r_i^q\ge z_i^q$ and coverage of $M$.

For the budget, every open bundle lies on the finite face of $\ell_{x_B}$.
Its nonnegative contributions sum to $\ell_{x_B}(x_B)=f(x_B)\le T$,
so \eqref{ext:eq:open-sum} holds.  At the fixed profile there are at most
$2^n$ capped vectors.  Lemma~\ref{ext:lem:budget} gives the $C_j(2^n)$ bound;
$M_i\le T$ gives the other bound.  All endpoint arithmetic above is finite.
\end{proof}

\begin{corollary}[Owner-count upper bound]\label{ext:cor:owner-upper}
Let $L_n=n\log2$ and $j_n=\max\{2,\lceil L_n/[\log(2+L_n)]^2\rceil\}$.
The rule with this report-independent parameter satisfies
\[
 \sum_iM_i(b)\le C_{j_n}(2^n)T
   =(\log2+o(1))\frac n{\log n}T,
\]
uniformly in dimension, technology and profiles.  For $n=1$ it is efficient.
\end{corollary}
\begin{proof}
Here $j_n\sim L_n/(\log L_n)^2$, $\alpha_{j_n}\sim\log n$, and
$j_n=o(n/\log n)$; substitute into $C_j$.
For $n=1$ and $b_1\ne0$, the off form is $\ell_0$, so its capped
contribution is $T$, and the endpoint marginal is also $T$.
Thus $r_1^q=M_1=T$; the zero bundle is immediate.
\end{proof}
\begin{remark}\label{ext:rem:direct}
The direct construction covers downward boundary discontinuities, for example
$f(x_1,x_2)=x_1$ when $x_2>0$ and zero otherwise.
It avoids limits of approximating technologies.  The canonical relative
support is a mathematical definition, not an efficient oracle; no regularity
or measurability of the resulting all-profile payment map is asserted.
\end{remark}

\subsection{Owner-count and dimension lower bounds}
\begin{definition}[Total-payment factors]\label{ext:def:factors}
Let $\mathcal Y_n(f)$ denote eligible rules on all of $(\R_+^d)^n$ and put
\[
 \mathfrak b_n(f)=\inf_{y\in\mathcal Y_n(f)}\sup_{b:T(b)>0}\frac{y(b)(N)}{T(b)}.
\]
The inner supremum is zero if $f\equiv0$.  These are total-payment factors,
not subsidy quantities.  Let $\mathfrak M_d$ be finite nonnegative min-linear
technologies on $\R_+^d$, and $\mathfrak F_d$ all finite monotone concave
homogeneous technologies there.  Define
\[
 B_n^{\rm own}=\sup_{d\ge1,f\in\mathfrak M_d}\mathfrak b_n(f),\quad
 B_n^{\rm conc}=\sup_{d\ge1,f\in\mathfrak F_d}\mathfrak b_n(f),\quad
 B_d^{\rm dim}=\sup_{n\ge1,f\in\mathfrak M_d}\mathfrak b_n(f).
\]
Owners, resources and technology are fixed inside each infimum; different
problems need not share a rule.  The first two suprema allow dimension to
grow; the last fixes dimension and permits arbitrarily many owners.
\end{definition}

\begin{theorem}[Owner-count range]\label{ext:thm:owner-range}
As $n\to\infty$ along all integers,
\[
 (1-o(1))\frac{\log n}{\log\log n}\le B_n^{\rm own}
 \le B_n^{\rm conc}\le\min\{n,C_{j_n}(2^n)\}
       =(\log2+o(1))\frac n{\log n}.
\]
The lower bound is a parameter translation of the weighted hierarchy,
not a new construction; it permits dimension to grow with $n$.
\end{theorem}
\begin{proof}
The upper follows from Theorem~\ref{ext:thm:general} and
$\mathfrak M_d\subseteq\mathfrak F_d$.
For the lower, use $f_k(x)=\min_{1\le e\le k}x_e$.
Restrict $km$ owners to axis-supported bundles, with $m$ owners per axis.
An eligible all-bundle rule restricts to a valid scalar-report multi-glove
rule, and its restricted total-payment factor cannot be larger.
Thus Lemma~\ref{lem:weighted-hierarchy} applies to $\mathfrak b_{km}(f_k)$.
If $n\ge n'$, restricting an $n$-owner rule to zero reports for the last
$n-n'$ owners and discarding their nonnegative payments yields an eligible
$n'$-owner rule.  All original coalition constraints and unilateral comparisons
remain, so $\mathfrak b_n(f)\ge\mathfrak b_{n'}(f)$ without population incentives.

For an integer $h\ge2$ set
\begin{align*}
 R&=\lceil\log(h+1)\rceil, &s&=\lceil\sqrt R\rceil, &m&=hs,\\
 \alpha&=1/s, &A&=(h+1)m-h^2, &\lambda&=As,\\
 Z&=\sum_{r=0}^h\frac{\lambda^r}{r!}, &K_h&=\lceil2RZ\rceil, &\nu_h&=K_hm.
\end{align*}
Here $s\ge2$, so $m>h$, and $K_hp_0=K_h/Z\ge2R\ge2\log(h+1)$.
Equation~\eqref{eq:weighted-bound}, with hierarchy height $h$, yields
\begin{equation}\label{ext:eq:owner-lower-finite}
 \mathfrak b_{\nu_h}(f_{K_h})
 \ge\frac{h+1-h/s}{(1+1/s)(1+h/(h+1)^2)}=(1-o(1))h.
\end{equation}
Also $\log\lambda=2\log h+O(\log\log h)$ and $\lambda\ge h$.
The terms of $Z$ increase up to index $h$, so
$\lambda^h/h!\le Z\le(h+1)\lambda^h/h!$.
Stirling's formula gives $\log Z=(1+o(1))h\log h$ and hence
$\log\nu_h=(1+o(1))h\log h$.
For fixed $0<\epsilon<1$ choose
$h=\lfloor(1-\epsilon)\log n/\log\log n\rfloor$.
Then $\log\nu_h\le(1-\epsilon+o(1))\log n<\log n$ for all large $n$.
Padding to $n$ owners and \eqref{ext:eq:owner-lower-finite} give
$B_n^{\rm own}\ge(1-\epsilon-o(1))\log n/\log\log n$.
Taking the liminf and then $\epsilon\downarrow0$ proves the all-integer lower.
\end{proof}

\subsection{Marginal floors for the original core}
At a profile define the original efficient core and its coordinate benchmarks
by $\core(b)=\{y:y(N)=T,\ y(S)\ge v_b(S)\ \forall S\}$ and
$H_i(b)=\max_{y\in\core(b)}y_i$.  Put
\begin{equation}\label{ext:eq:marginal}
 U_i=T-f(c-b_i),\qquad
 \rho_i=\max_{e:c_e>0}b_{ie}/c_e\quad(\rho_i=0\text{ if }c=0).
\end{equation}

\begin{lemma}[Marginal contributions]\label{ext:lem:marginal}
For concave $f$, the core is nonempty and compact,
$0\le H_i\le U_i\le T$, and $U$ is eligible.  Furthermore
\[
 \sum_iU_i\le\min\{d,n\}T,\qquad
 \sum_iU_i\le\min\{K,d,n\}T\quad\text{for a finite family of $K$ forms}.
\]
\end{lemma}
\begin{proof}
The grand relative support gives the finite allocation $y_i=\ell_c(b_i)$.
It is nonnegative, sums to $\ell_c(c)=T$, and covers each coalition by
majorization.  Singleton constraints make the core a closed subset of
$[0,T]^n$, hence compact.  Its complementary-coalition constraint gives
$y_i\le T-v_b(N\setminus\{i\})=U_i$.  Thus $U$ dominates a covering core
allocation.  Also $0\le U_i\le T$ and its subtracted term omits the whole
own bundle, proving own-report monotonicity.
Since $c-b_i\ge(1-\rho_i)c$, monotonicity and homogeneity imply
$U_i\le\rho_iT$.  Moreover
\[
 \sum_i\rho_i\le\sum_{e:c_e>0}\sum_i b_{ie}/c_e
   =|\operatorname{supp}(c)|\le d,\qquad \rho_i\le1.
\]
For finite $V$, group owners by a chosen off minimizer $w$ of $w\cdot(c-b_i)$.
For a nonempty group $G$, writing $A=w\cdot c\ge T$ gives
\[
 \sum_{i\in G}U_i=|G|(T-A)+w\cdot\sum_{i\in G}b_i
 \le|G|T-(|G|-1)A\le T.
\]
There are at most $\min\{K,n\}$ groups.  This last cancellation is used only
for ordinary finite forms, never for extended forms taking infinity at $c$.
\end{proof}

\begin{proposition}[Original-core floors]\label{ext:prop:floor}
For concave $f$, let a fixed $D\ge1$ satisfy $\sum_iU_i\le DT$ at every
profile.  One may use $D=d$, $D=n$, or $D=K$ for finite min-linear $f$.
For any eligible $y$, the rule $y_i'=y_i+U_i/D$ is eligible, protects
$y_i'\ge H_i/D$, and costs at most $y(N)+T$ in total.
With $D=K$ this preserves the sharp leading coefficient of
Section~\ref{prod:section}.  The benchmarks use the original core.
\end{proposition}
\begin{proof}
Adding nonnegative own-report monotone payments preserves eligibility.
The inequalities $U_i\ge H_i$ and $\sum_iU_i\le DT$ give the assertions.
One additional factor of $T$ is lower order than $T\log K/\log\log K$.
\end{proof}

\subsection{Resource dimension and technology-aware fallbacks}
\begin{theorem}[Dimension fallback without concavity]\label{ext:thm:dimension}
For every finite monotone homogeneous $f$, even without concavity,
$P_i=T\rho_i$ is eligible and
\[
 \sum_iP_i\le\min\{n,d_+(c)\}T\le dT,
 \qquad d_+(c)=|\operatorname{supp}(c)|.
\]
Thus $\mathfrak b_n(f)\le\min\{n,d\}$.
\end{theorem}
\begin{proof}
With other reports fixed, every ratio is $a/(a+A)$, where $A\ge0$ is fixed
and $a$ increases with the corresponding own coordinate; use $0/0=0$.
If $A=0$, the ratio jumps from zero to one and remains monotone.
Thus both $\rho_i$ and $T$ increase weakly with $b_i$, and their nonnegative
product does too.  For a coalition put $r_S=\sum_{i\in S}\rho_i$.
The coordinate inequalities $b_i\le\rho_ic$ imply
$v_b(S)\le f(r_Sc)=r_ST=\sum_{i\in S}P_i$.
If $r_S=0$ this uses $f(0)=0$; at $T=0$ all worths vanish.
The budget is $\sum_i\rho_i\le\min\{n,d_+(c)\}$, proved above without
concavity.  Zero aggregate coordinates contribute nothing.
\end{proof}
The maximum normalized coordinate share is familiar from multi-resource
allocation~\cite{Ghodsi}.  That literature's demand-report allocation model
and strategy-proofness claim differ from these supplier-side payments.

\begin{proposition}[Exact dimension factor without concavity]\label{ext:prop:sharp-d}
For every dimension $d$, the worst optimal total-payment factor over all
finite monotone homogeneous technologies and all owner counts is exactly $d$.
This does not establish a sharp dimension factor for concave technologies.
\end{proposition}
\begin{proof}
Take $f(x)=\max_ex_e$ and $d$ unit-basis owner bundles.
Then $T=1$ and every singleton has value one, forcing payment at least $d$.
The upper is Theorem~\ref{ext:thm:dimension}.
The example persists with strictly positive bundles: for $0<\epsilon\le1$,
let $b_i=e_i+\epsilon\sum_{e\ne i}e_e$.  Every singleton still has value one,
while $T=1+(d-1)\epsilon$, forcing factor at least
$d/[1+(d-1)\epsilon]$.  With $\epsilon=d^{-2}$ this is $(1-o(1))d$.
For $d\ge2$ this technology is convex and not concave, so the argument does
not apply to $\mathfrak M_d$ or $\mathfrak F_d$.
\end{proof}

\begin{theorem}[Concave dimension gap]\label{ext:thm:dim-gap}
As $d\to\infty$ along all integers,
\[
 (1-o(1))\frac{\log d}{\log\log d}\le B_d^{\rm dim}\le d.
\]
The lower already holds for $f(x)=\min_ex_e$ with axis-supported owners.
Whether an $o(d)$ dimension-only upper bound holds for $\mathfrak M_d$ or
$\mathfrak F_d$ is not determined here.
\end{theorem}
\begin{proof}
The upper is Theorem~\ref{ext:thm:dimension}.
Use the parameters $R,s,m,\alpha,A,\lambda,Z,K_h$ from
Theorem~\ref{ext:thm:owner-range}.  For every integer $D\ge K_h$, the weighted
hierarchy with $k=D$ pools still has $Dp_0=D/Z\ge2R$.
The same bound \eqref{ext:eq:owner-lower-finite} therefore yields
$\mathfrak b_{Dm}(f_D)\ge(1-o(1))h$; the owner count $Dm$ is allowed to grow
in $B_D^{\rm dim}$.  Since $\log K_h=(1+o(1))h\log h$, choosing
$h=\lfloor(1-\epsilon)\log D/\log\log D\rfloor$ gives $D\ge K_h$ for all
large $D$.  The resulting liminf is at least $1-\epsilon$;
let $\epsilon\downarrow0$.
\end{proof}

\begin{proposition}[Technology-aware refinements]\label{ext:prop:dual-share}
For finite min-linear $f$ define, with $0/0=0$,
\[
 \theta_i=\max_{w\in V}\frac{w\cdot b_i}{w\cdot c},\qquad Q_i=T\theta_i.
\]
The rule $Q$ is eligible, $\theta_i\le\rho_i$, and
$\sum_iQ_i\le\min\{n,d_+(c)\}T$.  For $|V|=1$ it is efficient.
If every $w\in V$ lies in the nonnegative cone of fixed vectors
$a_1,\ldots,a_r\in\R_+^d$, where $r\ge1$ is an integer, then
\[
 \theta_i\le\sigma_i:=\max_{1\le l\le r}\frac{a_l\cdot b_i}{a_l\cdot c},
 \qquad\sum_i\theta_i\le\sum_i\sigma_i\le r,
\]
so the budget improves to $\min\{n,d_+(c),r\}T$.
The marginal rule also satisfies $U_i\le P_i=T\rho_i$, with total
$\min\{K,d,n\}T$.
\end{proposition}
\begin{proof}
For fixed $w$, its ratio has the same $a/(a+A)$ monotonicity as above.
Taking the maximum and multiplying by monotone nonnegative $T$ preserves it.
For $T>0$, choose a form $w^*$ with $w^*\cdot c=T$.  Then
\[
 v_b(S)\le w^*\cdot\sum_{i\in S}b_i
     =T\sum_{i\in S}\frac{w^*\cdot b_i}{w^*\cdot c}
     \le\sum_{i\in S}Q_i.
\]
At $T=0$ all payments and worths vanish.
Each nonzero-denominator dual ratio is a weighted average of the coordinate
shares $b_{ie}/c_e$, using weights $w_ec_e$, so is at most $\rho_i$.
For one form its shares sum to one at $T>0$, proving efficiency.
If $w=\sum_l\mu_la_l$ with $\mu_l\ge0$, its ratio is similarly a weighted
average of generator ratios, with weights $\mu_l(a_l\cdot c)$.
Zero denominators also have zero numerators.  Finally
$\sum_i\sigma_i\le\sum_l\sum_i(a_l\cdot b_i)/(a_l\cdot c)\le r$.
The marginal claims are Lemma~\ref{ext:lem:marginal}.
No bound in terms of ordinary signed linear rank is asserted.
\end{proof}

\begin{remark}[Rule limitations and ex ante choice]\label{ext:rem:limits}
For $f(x)=\sum_ex_e$ and unit-basis owners, $P$ costs $dT$ whereas $Q$ is
efficient.  For $f(x)=\min_ex_e$ at that profile, $U$ costs $dT$ although
an efficient core allocation exists.  These are limitations of particular
rules, not lower bounds on all eligible rules.
Choosing among the valid rules in advance from fixed $n,d,K$ gives factor
$\min\{n,d,C_j(2^n)\}$ for general concave $f$, and
$\min\{n,d,C_j(\min\{K,2^n\})\}$ for finite min-linear $f$.
Coordinatewise minima or report-dependent switching are not asserted valid.
The convex-analysis, marginal/core, normalized-share, and Chernoff ingredients
are established tools; no historical priority, sharp owner/dimension rate,
or efficient evaluation for arbitrary $f$ is claimed.
\end{remark}
